\clearpage
\fancyhead[L]{\small LAPLACE. EXPOSITION DU SYSTÈME DU MONDE.}
\SourceLeaf{505}
\begin{center}\small Göttingische gelehrte Anzeigen. Stück 119. Seite 1185..1192. 1808. Juli 23.\\[3pt]\rule{.72\linewidth}{.3pt}\end{center}

\emph{Exposition du système du monde, par} M. \textsc{Laplace}, \emph{chancelier du Senat-Conservateur etc. Troisième édition, revue et augmentée par l'auteur. Chez Courcier,} 1808. 405 Seiten in Quart. Mit dem Bildnisse des Verfassers.

Der Gegenstand, der Plan, die Behandlungsart und der Werth dieses in seiner Art classischen Werks, wovon die frühern Ausgaben sich in Jedermanns Händen befinden, sind zu bekannt, als dass wir uns jetzt noch dabei aufhalten dürften; von dem Inhalte findet man auch schon im 97. u. 131. Stück dieser Blätter vom Jahr 1798, nach der Deutschen Uebersetzung der ersten Ausgabe, eine ausführliche Anzeige. Die zweite Ausgabe, welche 1799 erschien, hatte eben keine bedeutende Aenderungen erlitten. Wir begnügen uns also hier, nur die vornehmsten Aenderungen und Zusätze zu berühren, wodurch diese dritte Ausgabe sich von den frühern unterscheidet. Dass die Vermehrungen erheblich sein müssen, zeigt schon die grössere Seitenzahl, da die zweite Ausgabe nur 351 S. hatte. Diese Vermehrungen bestehen theils aus einigen grössern Zusätzen, theils aus einer grossen Menge kleinerer, wodurch der Inhalt der abgehandelten Materien grössere Vollständigkeit, und ihre Darstellung noch innigern Zusammenhang, noch mehr Evidenz, Fruchtbarkeit und Interesse erhält. Ein besonderer Vorzug dieses Werks besteht darin, dass von den numerischen Resultaten der Astronomie, in so fern sie in eine allgemeine Darstellung der Umrisse gehören können, die neuesten und bewährtesten Bestimmungen mitgetheilt werden, in welcher Rücksicht es um so mehr als Autorität gelten kann, da sehr viele davon durch die eigenen tiefen Untersuchungen des Verf. begründet oder veranlasst sind. Wir glauben Manchem einen angenehmen Dienst zu erweisen, wenn wir die vornehmsten davon, die seit der 2. Ausgabe neue Verbesserungen erhalten haben, in dieser Anzeige ausheben. 

Das \emph{erste Buch}, welches sich mit den \emph{scheinbaren} Bewegungen der Himmelskörper beschäftigt, hat zwei neue Kapitel erhalten, nemlich \emph{das zehnte} über die teleskopischen Planeten, Ceres, Pallas, Juno und Vesta, und \emph{das eilfte} über die Bewegung der Planeten um die Sonne, wogegen das zweite des zweiten Buchs mit der gleichen Ueberschrift weggeblieben ist. In der That gewinnt die Darstellung der Gründe für die wahre Weltordnung dadurch an Evidenz, wenn schon vorher ausgeführt ist, dass die Bewegungen der Planeten uns genau so erscheinen, als geschähen sie in Epicykeln, deren Mittelpunkt stets mit der Sonne zusammenfiele, und mit dieser sich um die Erde bewegte. Dies lässt sich bei Venus und Merkur aus ihren Phasen in Verbindung mit den successiven Veränderungen ihrer scheinbaren Durchmesser, bei Jupiter und Saturn aus den Verfinsterungen der Trabanten und den Verschwindungen des Ringes beweisen. Diese Folgerungen aus den Thatsachen gehören also allerdings in den ersten Abschnitt, da sie von der Wahl des Weltsystems unabhängig sind; aber vor Erfindung der Fernröhre waren diese Thatsachen freilich unbekannt, und \textsc{Ptolemaeus} konnte daher in seinem Systeme bei jedem Planeten nur das Verhältniss der Halbmesser der Cirkel und Epicykel angeben, und musste ihre absoluten Grössen unbestimmt lassen.

Von verbesserten Angaben numerischer Resultate in diesem Buche bemerken wir folgende, wobei wir uns auf die Seitenzahl der zweiten Ausgabe beziehen: dass dabei überall die neuen Decimaleintheilungen zum Grunde liegen, brauchen wir nicht zu erinnern. S. 6 Schiefe der Ekliptik für 1801, $26^{\circ}\,07315$, Säcularabnahme derselben (S. 11) für gegenwärtiges Jahrhundert $160''85$, S. 16 der siderische Tag 0.997269672 des mittlern Sonnentages. S. 17 das tropische Jahr jetzt 365.242264 Tage. S. 20 Sideralumlauf des Mondes 27.3216610716 Tage für den Anfang des 19. Jahrh., Sideralumlauf der Erdnähe 3232.58075
\SourceLeaf{506} Tage, des Knoten 6793.42118 Tage; die Epoche für 1801 (Mitternacht vor dem 1. Januar in Paris) für die mittlere Länge des Mondes $124^{\circ}01299$, für die Erdnähe $295^{\circ}66824$, für den aufsteigenden Knoten $17^{\circ}6933$. Ferner, Excentricität der Mondbahn 0.0548553, Neigung $5^{\circ}7222$, synodischer Umlauf des Mondes 29.58058817896 Tage. Die sich auf die Planeten beziehenden neuen Bestimmungen werden wir weiter unten zusammenstellen. S. 51 die grosse Axe der Nutationsellipse $59''56$, die kleine in Theilen des Parallelkreises $111''30$. S. 55 der nordische Grad unter der mittlern Breite $73^{\circ}7$, welcher nach den Franz. Astronomen 100696 Meter hatte, wird nach der neuen Schwed. Gradmessung nur 100316.1 Meter. Vielleicht würde auch der von \textsc{Lacaille} am Cap gemessene Grad, auf welchen man wohl etwas zu viel Gewicht gelegt zu haben scheint, um die Unregelmässigkeit der Figur der Erde zu beweisen, bei wiederholter Messung eine nicht unbeträchtliche Aenderung erleiden. Der Erzählung der vornehmsten Methoden, die Länge zu bestimmen, hat \textsc{Laplace} den Wunsch beigefügt, dass alle Nationen sich dahin vereinigen möchten, die Länge, anstatt von der vornehmsten Landessternwarte, von irgend einem physisch vorzüglich ausgezeichneten Punkte zu zählen, wozu sich besonders der Pik von Teneriffa gut eignen würde. S. 81 das Verhältniss des specifischen Gewichts der atmosphärischen Luft zu dem des Quecksilbers bei 0.76 Barometerstande und beim schmelzenden Eise wie 1 zu 10477.9. Der ganze Abschnitt über die Höhenmessungen mit dem Barometer und über die Strahlenbrechung ist mit mehr Ausführlichkeit behandelt.

Im \emph{zweiten Buche}, von den wahren Bewegungen der Himmelskörper, finden wir wenige Aenderungen, als dass im VI. (jetzt V.) Capitel die Schätzung der Wahrscheinlichkeit, dass zwei beobachtete Cometen von beinahe gleichen Elementen nur Einer sind, weggeblieben ist, und dagegen einige Vermuthungen über die Verwendung der Sonnenwärme auf den Cometen gewagt sind. Gänzlich umgearbeitet ist hingegen die Tafel der Bestimmungsstücke der sämmtlichen Planeten, die wir daher ihrer Wichtigkeit wegen hier ganz aufnehmen. Dass \textsc{Laplace} jetzt die Epoche von der Mitternacht zwischen dem 31. December und 1. Januar zählt, ist bereits erwähnt.

\begin{center} \Needspace{140pt}
{\small\setlength{\tabcolsep}{6pt}
\begin{tabular}{l|r|r}
 & \shortstack{Mittlere Länge 1801\\Pariser Meridian} & \shortstack{Siderischer Umlauf\\in Tagen} \\
\hline
Merkur & $182^{\circ}15647$ & $87.96925804$ \\
Venus & $11.93672$ & $224.70082399$ \\
Erde & $111.28179$ & $365.25638350$ \\
Mars & $71.24145$ & $686.9796186$ \\
Jupiter & $124.67781$ & $4332.5963076$ \\
Saturn & $150.38010$ & $10758.9698400$ \\
Uranus & $197.54244$ & $30688.7126872$ \\
\end{tabular}} \end{center}

\begin{center} \Needspace{140pt}
{\small\setlength{\tabcolsep}{6pt}
\begin{tabular}{l|r|r|r}
 & \shortstack{Halbe gross\\Axe} & \shortstack{Excentricität\\1801} & \shortstack{Säcularänderung} \\
\hline
Merkur & $0.3870981$ & $0.20551494$ & $+ \text{c}.000003867$ \\
Venus & $0.7233323$ & $0.00685298$ & $- 0.000062711$ \\
Erde & $1.0000000$ & $0.01685318$ & $- 0.00\text{c}041632$ \\
Mars & $1.5236935$ & $0.09313400$ & $+ 0.000090176$ \\
Jupiter & $5\;2027911$ & $0.04817840$ & $+ 0.000159350$ \\
Saturn & $9.5387705$ & $0.05616830$ & $- 0.000312402$ \\
Uranus & $19.1833050$ & $0.04667030$ & $- 0.000025072$ \\
\end{tabular}} \end{center}

\SourceLeaf{507}
\begin{center} \Needspace{140pt}
{\small\setlength{\tabcolsep}{6pt}
\begin{tabular}{l|r|r}
 & \shortstack{Sonnennähe 1801} & \shortstack{Siderische\\Säcularbewegung} \\
\hline
Merkur & $82^{\circ}6256$ & $+ 1801''10$ \\
Venus & $142.9077$ & $- 826.63$ \\
Erde & $110.5571$ & $+ 3641.40$ \\
Mars & $369.3407$ & $+ 4884.05$ \\
Jupiter & $12.3812$ & $+ 2048.95$ \\
Saturn & $99.0549$ & $+ 5978.60$ \\
Uranus & $185.9574$ & $+ 738.69$ \\
\end{tabular}} \end{center}

\begin{center} \Needspace{140pt}
{\small\setlength{\tabcolsep}{6pt}
\begin{tabular}{l|r|r}
 & \shortstack{Neigung der Bahn\\1801} & \shortstack{Säcularänderung} \\
\hline
Merkur & $7^{\circ}78058$ & $+ 56''12$ \\
Venus & $3.76936$ & $- 14.05$ \\
Mars & $2.05663$ & $- 0.47$ \\
Jupiter & $1.46034$ & $- 69.78$ \\
Saturn & $2.77102$ & $- 47.88$ \\
Uranus & $0.85990$ & $+ 9.67$ \\
\end{tabular}} \end{center}

\begin{center} \Needspace{140pt}
{\small\setlength{\tabcolsep}{6pt}
\begin{tabular}{l|r|r}
 & \shortstack{Aufsteig. Knoten\\1801} & \shortstack{Siderische\\Säcularbewegung} \\
\hline
Merkur & $51^{\circ}0651$ & $- 2414''41$ \\
Venus & $83.1972$ & $- 5770.99$ \\
Mars & $53.3605$ & $- 7186.65$ \\
Jupiter & $109.3624$ & $- 4869.04$ \\
Saturn & $124.3662$ & $- 6995.25$ \\
Uranus & $10.9488$ & $- 11104.81$ \\
\end{tabular}} \end{center}

Die Elemente der vier neuen Planeten, welche hier gegeben werden, sind bei der Ceres die 11., bei der Pallas die 9., bei der Juno die 6. von \textsc{Gauss}, bei der Vesta diejenigen, welche Hr. \textsc{Burkhardt} in der Connoissance des tems 1809 gegeben hat, und bei denen nur erst wenige Beobachtungen benutzt waren. Wir lassen sie daher hier weg, da in Deutschland längst genauere bekannt sind (Monatl. Corresp. Febr. 1808; Gött. gel. Anz. 1808 St. 14. 40. 107).

Auch das \emph{dritte Buch}, über die Gesetze der Bewegung, hat nur ein paar kleine Zusätze erhalten, die sich auf die Darstellung der ersten Grundsätze der Dynamik in Gleichungen, auf das Princip der kleinsten Wirkung, auf den Begriff \emph{Masse}, und die beiden Arten des Gleichgewichts beziehen.

Im \emph{vierten Buche}, über die Theorie der allgemeinen Schwere, ist das Kapitel über die Störungen der elliptischen Bewegungen der Planeten das zweite geworden; zwei neue Kapitel, über die Trabanten des Saturn und Uranus, und über diejenige Anziehung der kleinsten Theile der Körper, welche nur in unmerklichen Entfernungen merklich ist (attraction moléculaire), sind hinzugekommen. Dieses letztere Kapitel gibt eine Uebersicht dieser Theorie, die \textsc{Laplace} in zweien Supplementen zu seiner Mécanique céleste (De l'action capillaire, und Supplément à la théorie de l'action capillaire) entwickelt, und mit so überraschend glücklichem Erfolge zu einer mathematischen Erklärung der Strahlenbrechung, der Phänomene der Haarröhren, des Anziehens und Abstossens kleiner, auf einer Flüssigkeit schwimmenden, Körper, des Zusammenhängens einer Scheibe mit einer Flüssigkeit, des Schwimmens kleiner fester Körper in einer specifisch leichtern Flüssigkeit u.s.w. angewandt hat. Diese Entdeckungen machen Epoche in der Physik: sie brechen die Bahn zu einer Wissenschaft, die für die Natur in Beziehung auf die kleinsten Theile der Körper das sein wird, was die Gravitationslehre für die Natur im Grossen ist. Eine ausführlichere Anzeige davon müssen wir uns aber auf eine andere Gelegenheit versparen.

\SourceLeaf{508}
Von kleinern Abänderungen und Zusätzen bemerken wir aus diesem Abschnitt noch folgende: Die grosse Gleichung des Saturn steigt auf $9111''41$, und ihre Periode ist von $921\frac{1}{2}$ Jahren; die grosse Gleichung des Jupiter ist $3720''36$. Das Verhältniss der Massen der Planeten zu der Masse der Sonne ist wie 1 zu den Zahlen 2025810 bei Merkur; 356632 bei Venus; 337086 bei der Erde; 2546320 bei Mars; 1067.09 bei Jupiter; 3534.08 bei Saturn; 19504 bei Uranus. Der scheinbare Jupitersdurchmesser in der Distanz 1 ist $599''151$; das Verhältniss der Schwere eines Körpers unter dem Erdäquator, unter dem Jupitersäquator, und unter dem Sonnenäquator wie die Zahlen 1000, 2566, 27933. Bei der Theorie des Mondes sind mehrere Zusätze eingeschaltet, über die Ungleichheiten, welche von der Abplattung der Erde und von der Parallaxe der Sonne abhängen, ferner über die Säcularungleichheiten in der Bewegung des Knoten und der Apsiden, und über die vor einigen Jahren entdeckte Gleichung, deren Periode von 184 Jahren ist. --- Auch der Abschnitt über die Jupiterstrabanten hat verschiedene weitere Ausführungen erhalten.

Das \emph{fünfte Buch} enthält die Geschichte der Astronomie: auch hier nur Umrisse, die aber, von einer solchen Hand gezeichnet, für den Kenner wie für den Liebhaber ein hohes Interesse haben. Le tableau des progrès de la plus sublime des sciences naturelles, sagt \textsc{Laplace}, toujours croissans au milieu même des révolutions des empires, pourra consoler des malheurs dont les récits remplissent les annales de tous les peuples. Auch dieser Theil des Werks ist mit manchen Zusätzen und interessanten Reflexionen bereichert. Dahin gehören die Bemerkungen über die Spuren von astronomischen Kenntnissen, die man bei den Eingebornen von Mexico und Peru antraf. Jene hatten eine sehr genaue Kenntniss von der Länge des tropischen Jahrs; sie bedienten sich einer Einschaltungsmethode, wobei dasselbe zu $365\frac{25}{104}$ Tagen oder 365 Tagen 5 Stunden 46 Min. 9 Sec. vorausgesetzt wird. Dabei ist es sehr merkwürdig, dass ihnen die Zeitabtheilung in Wochen, welche man bei allen Völkern der alten Welt findet, unbekannt war, und dass sie statt derselben eine Periode von fünf Tagen hatten.

Am Schlusse des Werks finden sich noch sechs historische Anmerkungen. Die erste bezieht sich auf ein paar Chinesische Beobachtungen von \textsc{Tcheou-Kong}, wonach ungefähr um das Jahr 1100 vor unsrer Zeitrechnung die Schiefe der Ekliptik $26^{\circ}5563$, und die gerade Aufsteigung des Sterns $\varepsilon$ im Wassermann $297^{\circ}8096$ war; nach \textsc{Laplace}'s Formeln sollte für jenes Jahr jene $26^{\circ}5161$, und diese $298^{\circ}7265$ sein; um die letztern Zahlen in völlige Uebereinstimmung zu bringen, brauchte man nur noch 54 Jahre weiter zurück zu gehen. Die zweite betrifft eine Nachricht, die uns \textsc{Geminus} über die Kenntnisse der Chaldäer vom Mondslaufe aufbehalten hat. Die dritte Note betrifft \textsc{Pytheas} Beobachtung des Solstitiums zu Marseille; die vierte vergleicht \textsc{Hipparchs} Angaben für die Bewegung des Mondes in Beziehung auf die Sonne, die Erdnähe und den Knoten mit den Bestimmungen, die aus \textsc{Laplace}'s Theorie der Säcularungleichheiten folgen; die fünfte vergleicht die grösste Mittelpunktsgleichung der Sonne, die Lage der Sonnenferne, die Länge des Jahrs und die Schiefe der Ekliptik nach den Bestimmungen der Araber, mit den neuesten Angaben; endlich die sechste stellt die Bestimmungen der Schiefe der Ekliptik von \textsc{Tcheou-Kong} 1100 J. vor Chr. Geb., von \textsc{Pytheas} 350 vor Chr. Geb., von \textsc{Ibn Junis} im J. 1000, von \textsc{Cocheou-Kong} im Jahr 1280, von \textsc{Ulug-Beigh} im Jahr 1437, und die neuesten von 1801 mit den Resultaten der Theorie zusammen.
